\documentclass[9pt, mathserif, professionalfonts]{beamer}


\usepackage[ngerman]{babel}

\usepackage{fontspec}

\setsansfont{Tex Gyre Heros}


\setlength{\parindent}{0pt}
\usepackage[babel,german=quotes]{csquotes}
\usepackage{eurosym}

\usepackage{amsmath}
\usepackage{mathtools}
\usepackage{amssymb}
\usepackage{xfrac}
\usepackage{unicode-math}
\setmathfont{Tex Gyre Pagella Math}

\usepackage[locale=DE, separate-uncertainty=true, per-mode=fraction]{siunitx}

\usepackage{booktabs}
\usepackage{tabulary}
\usepackage{threeparttable}

\usepackage[]{graphicx}
\usepackage[labelfont=bf]{caption}
\setbeamertemplate{caption}[numbered]

\usepackage[]{hyperref}
\usepackage{braket}
\usepackage{xcolor}
\usepackage{enumerate}

\usepackage{listings}
\usepackage{lstautogobble}

\lstnewenvironment{lstverbatim}[1][]{
  \lstset
  {
    autogobble,
    language={[LaTeX]TeX},                
    basicstyle=\ttfamily,
    texcsstyle=*\color{TUdarkgreen},
    keywordstyle=\color{TUorange},
    commentstyle=\color{blue},
    morekeywords = {hallo},
    otherkeywords={$,$, &, \{, \}, \[, \],  ^, _},
    moretexcs={text, num, SI, si, kilo, gram, befehl, subject, titlehead, publishers,
               metre, per, second, squared, micro, subtitle, maketitle, tableofcontents,
               subsection, subsubsection, part, chapter, paragraph, subparagraph,
               iint, iiint},
    columns=fullflexible,
    breaklines=true,
    keepspaces=true,
    aboveskip=-0.2em,
    belowskip=-0.2em
  }
}{}


\title{Verfassen wissenschaftlicher Texte mit \LaTeX}
\author{PeP et Al. toolbox}
\date{2014}
\institute[Pep et Al. e.V.]{\includegraphics[width=0.2\textwidth]{peplogo.pdf}}
\titlegraphic{./latex-vs-word.pdf}

\usetheme{TUDo}

\begin{document}

\begin{frame}
\setcounter{framenumber}{0}
    \titlepage
\end{frame}

\begin{frame}
    \frametitle{Übersicht}
    \tableofcontents
\end{frame}

\section{Grundlagen}
\begin{frame}
    \frametitle{Warum \LaTeX \ benutzen?}
    \begin{itemize}
        \item hervorragender Text- und Formelsatz
        \item automatisierte Erstellung von Inhalts- und Literaturverzeichnis
        \item Tex-Dateien sind reine Text-Dateien \\
              $\Rightarrow$ kleine Datein, gut für Versionskontrolle geeignet
        \item sehr gute Vorlagen für wissenschaftliche Arbeiten 
        \item aber auch: Notensatz, Präsentationen 
        \item kostenlos, OpenSource, ausgezeichnete Dokumentionen
        \item erweiterbar durch zahlreiche und mächtige Pakete
        \item aus allen geläufigen Betriebssystem verfügbar
        \item Ausgabe nach PDF
    \end{itemize}
\end{frame}

\begin{frame}[fragile]
    \frametitle{Das Dokument}
    Diese drei Zeilen braucht jedes \LaTeX-Dokument:
    \begin{columns}[T]
        \begin{column}{0.4\textwidth}
            \begin{block}{\LaTeX-Code}
                \begin{lstverbatim}
                \documentclass{scrartcl}
                    % Präambel
                    .
                    .
                \begin{document}
                    % Inhalt des Dokuments
                    .
                    .
                \end{document}
                \end{lstverbatim}
            \end{block}
        \end{column}
        \begin{column}{0.55\textwidth}
            \begin{block}{documentclass}
                Hier wird die Dokumentenklasse definiert. \\
                Es folgt die \alert{Präambel} des Dokuments. 
                Hier werden globale Optionen gesetzt und zusätzliche Pakete eingebunden.\\
            \end{block}

            \vspace{0.55cm}

            \begin{block}{document-Umgebung}
                Hier wird das eigentliche Dokument erstellt.
            \end{block}
        \end{column}
    \end{columns}
\end{frame}

\begin{frame}[fragile]
    \frametitle{Syntax: Befehle}
    \LaTeX-Befehle beginnen stets mit einem Back-Slash.

    Obligatorische Argumente stehen in \{ \}, optionale Argumente stehen in [ ].
    \begin{columns}[T]
        \begin{column}{0.47\textwidth}
            \begin{block}{\LaTeX-Code}
                \begin{lstverbatim}
                \befehl[optional]{obligatorisch}
                \documentclass[paper=a4]
                              {scrartcl}
                \frac{1}{2}
                \end{lstverbatim}
            \end{block}
        \end{column}
        \begin{column}{0.47\textwidth}
            \begin{block}{Erklärung}
                Beispiel \\
                Setzt die Dokumentenklasse auf \emph{scrartcl} und das Papierformat auf DIN\,A4. \\
                Es gibt auch Befehle mit zwei oder mehr Pflichtargumenten, z.B. der Bruch.
            \end{block}
        \end{column}
    \end{columns}
\end{frame}

\begin{frame}[fragile]
    \frametitle{Standardpakete}
    Die hier aufgezählten Pakete sollten immer geladen werden, da sie wesentliche Funktionen bieten und wichtige Einstellungen vornehmen.
    \begin{columns}[T]
        \begin{column}{0.47\textwidth}
            \begin{block}{Paket}
                \begin{lstverbatim}
                \usepackage[ngerman]{babel}

                \usepackage{fontspec}
                \end{lstverbatim}
            \end{block}
        \end{column}
        \begin{column}{0.47\textwidth}
            \begin{block}{Funktion}
                Deutsche Spracheinstellungen für das Dokument. \\
                Schriftpaket. Wichtig, damit Umlaute korrekt dargestellt werden.
                Erlaubt einfaches ändern der verschiedenen Schriftarten.
            \end{block}
        \end{column}
    \end{columns}
\end{frame}

\section{Text erstellen}

\begin{frame}[fragile]{Text schreiben}
  \begin{block}{Beispiel}
    \begin{lstlisting}
      % Präambel
      \begin{document}
        Hallo, Welt!

        Dies ist ein dummer Beispieltext.
        Er soll zeigen, dass \LaTeX sich nicht um
        Zeilenumbrüche im Code    oder    zuviele
        Leerzeichen kümmert.

        Ein Absatz wird mit einer leeren Code-Zeile
        markiert.
      \end{document}
    \end{lstlisting}
  \end{block}
\end{frame}

\begin{frame}[fragile]{Konventionen für Text}
  \begin{itemize}
    \item Höchstens ein Satz pro Code-Zeile
    \item Absätze werden durch eine Leerzeile markiert
    \item Im Fließtext sollten keine Umbrüche mit \lstinline+\\+ erzwungen werden
  \end{itemize}
  \begin{alertblock}{Sonderzeichen}
    Viele Sonderzeichen sind \LaTeX-Steuerzeichen.
    Damit diese im Text genutzt werden können, muss meist ein \verb+\+ vorangestellt oder ein Befehl genutzt werden.
  \end{alertblock}
  \begin{CodeExample}{0.7}
    \begin{lstlisting}
      \# \$ \% \& \_ \{ \}
      \textbackslash \textasciicircum \textasciitilde
    \end{lstlisting}
  \CodeResult
    \strut
    \# \$ \% \& \_ \{ \} \\
    \textbackslash\ \textasciicircum\ \textasciitilde
  \end{CodeExample}
\end{frame}

\begin{frame}[fragile]{Textauszeichnung}
  Änderungen der Schrifteigenschaften sind mit diesen Befehlen möglich:
  \begin{CodeExample}{0.60}
    \begin{lstlisting}
      \textit{kursiv} \emph{kursiv}
      \textbf{fett}
      \textbf{\textit{fett-kursiv}}
      \textrm{Serifen-Schrift}
      \texttt{Mono-Schrift}
      \textsf{Sans-Serif-Schrift}
      \textsc{Kapitälchen}
    \end{lstlisting}
  \CodeResult
    \strut
    \textit{kursiv} \emph{kursiv} \\
    \textbf{fett} \\
    \textbf{\textit{fett-kursiv}} \\
    \textrm{Serifen-Schrift} \\
    \texttt{Mono-Schrift} \\
    \textsf{Sans-Serif-Schrift} \\
    \textsc{Kapitälchen}
  \end{CodeExample}

  \vspace{1em}
  Diese Befehle sollten sehr selten benutzt werden, semantischer Markup ist besser.
\end{frame}

\begin{frame}[fragile]{Schriftgrößen}
  Gelten immer für den aktuellen Block, z.\,B. in einer Umgebung oder zwischen \lstinline+{ }+
  \begin{CodeExample}{0.48}
    \begin{lstlisting}
      {\tiny tiny}
      {\small small}
      {\normalsize normal}
      {\large large}
      {\huge huge}
    \end{lstlisting}
  \CodeResult
    \begin{minipage}[c][5\baselineskip][c]{\textwidth}
      {\tiny tiny}
      {\small small}
      {\normalsize normal}
      {\large large}
      {\huge huge}
    \end{minipage}
  \end{CodeExample}
  \vspace{1em}
  \begin{block}{Alle Größen}
    \begin{lstlisting}
      \tiny, \scriptsize, \footnotesize, \small, \normalsize, \large, \Large, \LARGE, \huge, \Huge
    \end{lstlisting}
  \end{block}
  Auch diese Befehle sollten nur über semantischen Markup benutzt werden.
\end{frame}

\begin{frame}[fragile]{Inhalt auslagern}
  \begin{block}{Code}
    \begin{lstlisting}
      \documentclass[
  bibliography=totoc,     % Literatur im Inhaltsverzeichnis
  captions=tableheading,  % Tabellenüberschriften
  titlepage=firstiscover, % Titelseite ist Deckblatt
]{scrartcl}

% LaTeX2e korrigieren.
\usepackage{fixltx2e}
% Warnung, falls nochmal kompiliert werden muss
\usepackage[aux]{rerunfilecheck}

% deutsche Spracheinstellungen
\usepackage{polyglossia}
\setmainlanguage{german}

% unverzichtbare Mathe-Befehle
\usepackage{amsmath}
% viele Mathe-Symbole
\usepackage{amssymb}
% Erweiterungen für amsmath
\usepackage{mathtools}

% Fonteinstellungen
\usepackage{fontspec}
\defaultfontfeatures{Ligatures=TeX}

\usepackage[
  math-style=ISO,    % \
  bold-style=ISO,    % |
  sans-style=italic, % | ISO-Standard folgen
  nabla=upright,     % |
  partial=upright,   % /
]{unicode-math}

\setmathfont{Latin Modern Math}
\setmathfont[range={\mathscr, \mathbfscr}]{XITS Math}
\setmathfont[range=\coloneq]{XITS Math}
\setmathfont[range=\propto]{XITS Math}
% make bar horizontal, use \hslash for slashed h
\let\hbar\relax
\DeclareMathSymbol{\hbar}{\mathord}{AMSb}{"7E}
\DeclareMathSymbol{ℏ}{\mathord}{AMSb}{"7E}

% richtige Anführungszeichen
\usepackage[autostyle]{csquotes}

% Zahlen und Einheiten
\usepackage[
  locale=DE,                   % deutsche Einstellungen
  separate-uncertainty=true,   % Immer Fehler mit \pm
  per-mode=symbol-or-fraction, % m/s im Text, sonst Brüche
]{siunitx}

% chemische Formeln
\usepackage[version=3]{mhchem}

% schöne Brüche im Text
\usepackage{xfrac}

% Floats innerhalb einer Section halten
\usepackage[section, below]{placeins}
% Captions schöner machen.
\usepackage[
  labelfont=bf,        % Tabelle x: Abbildung y: ist jetzt fett
  font=small,          % Schrift etwas kleiner als Dokument
  width=0.9\textwidth, % maximale Breite einer Caption schmaler
]{caption}
% subfigure, subtable, subref
\usepackage{subcaption}

% Grafiken können eingebunden werden
\usepackage{graphicx}
% größere Variation von Dateinamen möglich
\usepackage{grffile}

% Standardplatzierung für Floats einstellen
\usepackage{float}
\floatplacement{figure}{htbp}
\floatplacement{table}{htbp}

% schöne Tabellen
\usepackage{booktabs}

% Seite drehen für breite Tabellen
\usepackage{pdflscape}

% Literaturverzeichnis
\usepackage{biblatex}
% Quellendatenbank
\addbibresource{lit.bib}
\addbibresource{programme.bib}

% Hyperlinks im Dokument
\usepackage[
  unicode,
  pdfusetitle,    % Titel, Autoren und Datum als PDF-Attribute
  pdfcreator={},  % PDF-Attribute säubern
  pdfproducer={}, % "
]{hyperref}
% erweiterte Bookmarks im PDF
\usepackage{bookmark}

% Trennung von Wörtern mit Strichen
\usepackage[shortcuts]{extdash}

\author{
  AUTOR A%
  \texorpdfstring{
    \\
    \href{mailto:authorA@udo.edu}{authorA@udo.edu}
  }{}%
  \texorpdfstring{\and}{, }
  AUTOR B
  \texorpdfstring{
    \\
    \href{mailto:authorB@udo.edu}{authorB@udo.edu}
  }{}
}
\publishers{TU Dortmund – Fakultät Physik}

      \begin{document}
        \input{Teil1.tex}
        \input{Teil2.tex}
        % .
      \end{document}
    \end{lstlisting}
  \end{block}
  \begin{itemize}
    \item Verschachtelung möglich
    \item Zur Aufteilung größerer Dokumente (z.B. diese Präsentation)
    \item Für häufig wiederverwendeten Code (Header, Erläuterungen zu Fehlerrechnung, …)
    \item Für per Skript erzeugte Tabelleninhalte
  \end{itemize}
\end{frame}

\begin{frame}[fragile]{
  Anführungszeichen
  \hfill\doc{http://mirrors.ctan.org/macros/latex/contrib/csquotes/csquotes.pdf}{csquotes}
}
  Die richtigen Anführungszeichen, wo die Satzzeichen hingehören und vieles mehr hängt von der Sprache ab.
  So macht man es richtig:
  \begin{Packages}
    \begin{lstlisting}
      \usepackage[autostyle]{csquotes}    % nach polyglossia
      \setotherlanguages{english, french} % andere Sprachen laden
    \end{lstlisting}
  \end{Packages}
  \begin{CodeExample}{0.60}
    \begin{lstlisting}
      foo \enquote{bar} baz
      \enquote{foo \enquote{bar} baz}
      \textenglish{\enquote{foo}}
      \textfrench{\enquote{foo}}
      \textcquote{root}{foo}
    \end{lstlisting}
  \CodeResult
    \strut
    foo \enquote{bar} baz \\
    \enquote{foo \enquote{bar} baz} \\
    \textenglish{\enquote{foo}} \\
    \textfrench{ \enquote{foo}} \\
    \textcquote{root}{foo}
    % work around polyglossia bug #68
    \directlua{polyglossia.desactivate_frpt()}
  \end{CodeExample}
\end{frame}

\section{Zahlen und Einheiten}

\begin{frame}[fragile]{Zahlen und Einheiten}
  \begin{itemize}
    \item Regeln zur Benutzung der SI-Einheiten: \\
      \url{http://www.bipm.org/utils/common/pdf/si_brochure_8_en.pdf}
    \item Einheiten werden aufrecht gesetzt
    \item Zwischen Zahl und Einheit steht ein kleines Leerzeichen
    \item Ab 5 Stellen wird ein kleines Leerzeichen als 1000er Trennzeichen genutzt:
  \end{itemize}

  \begin{CodeExample}{0.7}[Zahl mit Einheit]
    \begin{lstlisting}
      $5\,\mathrm{kg}$
    \end{lstlisting}
    \CodeResult
    $5\,\mathrm{kg}$
  \end{CodeExample}
  \begin{CodeExample}{0.7}[Zahl mit mehr als vier Stellen]
    \begin{lstlisting}
      $10\,000$
    \end{lstlisting}
    \CodeResult
    $10\,000$
  \end{CodeExample}
  \begin{CodeExample}{0.7}[Zehnerpotenz mit Unsicherheit in Klammern]
    \begin{lstlisting}
      $(5{,}34 \pm 0{,}54) \cdot 10^{-3}\,\mathrm{GeV}$
    \end{lstlisting}
    \CodeResult
    $(5{,}34 \pm 0{,}54) \cdot 10^{-3}\,\mathrm{GeV}$
  \end{CodeExample}
  
  \onslide<2->{%
  \begin{center}
    \LARGE  Das muss einfacher gehen
  \end{center}
  }
\end{frame}

\begin{frame}[fragile]{
  Das \texttt{siunitx}-Paket
  \hfill
  \doc{http://mirrors.ctan.org/macros/latex/contrib/siunitx/siunitx.pdf}{siunitx}
}
\begin{itemize}
    \item \texttt{siunitx} stellt Befehle zur Verfügung, die das korrekte Setzen von Zahlen und Einheiten stark vereinfachen
    \item Funktioniert in Fließtext und Matheumgebung
    \item[$\color{vertexDarkRed}\Rightarrow$] Dieses Paket sollte \alert{immer} und für \alert{jede} Zahl mit oder ohne Einheit verwendet werden.
\end{itemize}
  \begin{Packages}
    \begin{lstlisting}
      \usepackage[
        locale=DE,
        separate-uncertainty=true,  % Immer Fehler mit ±
        per-mode=symbol-or-fraction, % m/s im Text, sonst \frac
        % alternativ:
        % per-mode=reciprocal,      % m s^{-1}
        % output-decimal-marker=.,         % . statt , für Dezimalzahlen
      ]{siunitx}
    \end{lstlisting}
  \end{Packages}
\end{frame}

\begin{frame}[fragile]{\texttt{siunitx}: Zahlen mit \lstinline+\\num+}
  \begin{CodeExample}{0.7}[Zahlen mit automatischen 3er-Gruppen]
    \begin{lstlisting}
      \num{1.23456}
      \num{987654321}
    \end{lstlisting}
  \CodeResult
    \strut
    \num{1.23456} \\
    \num{987654321}
  \end{CodeExample}
  \begin{CodeExample}{0.7}[Einfaches Eingeben von 10er Potenzen]
    \begin{lstlisting}
      \num{6.022e23}
    \end{lstlisting}
  \CodeResult
    \strut
    \num{6.022e23}
  \end{CodeExample}
  \begin{CodeExample}{0.7}[Angabe von Fehlern]
    \begin{lstlisting}
      \num{1.54 +- 0.1}
      \num{1.54(10)}
      \num{1.54 \pm 0.1}
      \num[separate-uncertainty=false]{1.54 +- 0.1}
      \num{3.5(1)e6}
    \end{lstlisting}
  \CodeResult
    \strut
    \num{1.54 +- 0.1} \\
    \num{1.54(10)} \\
    \num{1.54 \pm 0.1} \\
    \num[separate-uncertainty=false]{1.54 +- 0.1} \\
    \num{3.5(1)e6}
  \end{CodeExample}
\end{frame}

\begin{frame}[fragile]{\texttt{siunitx}: Einheiten mit \lstinline+\\si+}
  \begin{CodeExample}{0.8}[Einheiten]
    \begin{lstlisting}
      \si{\meter\per\second}
      \si[per-mode=fraction]{\meter\per\second}
      \si{\meter\per\second\squared}
      \si[per-mode=reciprocal]{\gram\per\cubic\centi\meter}
      \si{\kelvin\tothe{4}}
    \end{lstlisting}
  \CodeResult
    \strut
    \si{\meter\per\second} \\
    \si[per-mode=fraction]{\meter\per\second} \\
    \si{\meter\per\second\squared} \\
    \si[per-mode=reciprocal]{\gram\per\cubic\centi\meter} \\
    \si{\kelvin\tothe{4}}
  \end{CodeExample}
  \begin{CodeExample}{0.80}[\texttt{per-mode=symbol-or-fraction}]
    \begin{lstlisting}
      \begin{equation}
        \si{\kilo\gram\meter\per\second\squared}
      \end{equation}
      $\si{\kilo\gram\meter\per\second\squared}$
    \end{lstlisting}
  \CodeResult
    \removedisplayskip
    \begin{minipage}[c][3\baselineskip][c]{\textwidth}
      \begin{equation}
        \si{\kilo\gram\meter\per\second\squared}
      \end{equation}
    \end{minipage}
    $\si{\kilo\gram\meter\per\second\squared}$
  \end{CodeExample}
  \begin{CodeExample}{0.80}[Meter mal Sekunde oder Millisekunde?]
    \begin{lstlisting}
      \si{\milli\second}
      \si{\meter\second}
      \si[inter-unit-product=\cdot]{\meter\second}
    \end{lstlisting}
  \CodeResult
    \strut
    \si{\milli\second} \\
    \si{\meter\second} \\
    \si[inter-unit-product=\cdot]{\meter\second}
  \end{CodeExample}
\end{frame}

\begin{frame}[fragile]{\texttt{siunitx}: Zahl mit Einheit: \lstinline+\\SI+}
  \begin{CodeExample}{0.6}[\lstinline+\\SI+ {=} Kombination aus \lstinline+\\num+ und \lstinline+\\si+]
    \begin{lstlisting}
      \SI{5}{\percent}
      \SI{10}{\celsius}
      \SI{2.5(1)e6}{\kilo\gram\square\meter\per\second\squared}
    \end{lstlisting}
  \CodeResult
    \strut
    \SI{5}{\percent} \\
    \SI{10}{\celsius} \\
    \SI{2.5(1)e6}{\kilo\gram\square\meter\per\second\squared}
  \end{CodeExample}
  \begin{description}
    \item[1. Argument] Kann alles, was \lstinline+\num+ kann
    \item[2. Argument] Kann alles, was \lstinline+\si+ kann
  \end{description}
  \begin{CodeExample}{0.6}[Winkel]
    \begin{lstlisting}
      \ang{5;55;59}
    \end{lstlisting}
  \CodeResult
    \strut
    \ang{5;55;59}
  \end{CodeExample}
\end{frame}

\section{Formelsatz}
\subsection{im Fließtext}
\begin{frame}[fragile]
   \frametitle{\$ \dots \$-Umgebung}
   Für mathematische Symbole, Variablen und kleine Formeln im Fließtext. 

    \begin{columns}[T]
        \begin{column}{0.6\textwidth}
            \begin{block}{\LaTeX-Code}
                \begin{lstverbatim}
                $x$ 
                $x^i$
                $x^12$ bzw. $x^{12}$ % Vorsicht
                $x_\mathrm{max}$
                $U(t) = U_0 \cdot \cos(\omega t)$
                \end{lstverbatim}
            \end{block}
        \end{column}
        \begin{column}{0.35\textwidth}
            \begin{block}{Ergebnis}
                $x$  \\
                $x^i$ \\
                $x^12$ bzw. $x^{12}$ \\ 
                $x_\mathrm{max}$ \\
                $U(t) = U_0 \cdot \cos(\omega t)$
            \end{block}
        \end{column}
    \end{columns}
\end{frame}
\subsection{Mathe-Umgebungen}
\begin{frame}[fragile]
    \frametitle{Die \texttt{equation}-Umgebung 
        \hfill\doc{http://ftp.fau.de/ctan/macros/latex/required/amslatex/math/amsldoc.pdf}{amsmath}
        \doc{http://ftp.uni-erlangen.de/mirrors/CTAN/macros/latex/contrib/mathtools/mathtools.pdf}{mathtools}
    }
   \begin{block}{benötigte Pakete}
    \begin{lstverbatim}
        \usepackage{amsmath}
        \usepackage{amssymb}
        \usepackage{mathtools}
    \end{lstverbatim}
   \end{block}
    Für eine abgesetzte Gleichung, die automatisch nummeriert wird.
    \begin{columns}[T]
        \begin{column}{0.6\textwidth}
            \begin{block}{\LaTeX-Code}
                \begin{lstverbatim}
                \begin{equation}
                    \nabla \cdot \vec{E} = 
                    \frac{\rho} {\epsilon_0}
                    \label{eq:maxwell1}
                \end{equation}
                \end{lstverbatim}
            \end{block}
        \end{column}
        \begin{column}{0.35\textwidth}
            \begin{block}{Ergebnis}
                \begin{equation}
                    \nabla \cdot \vec{E} = 
                    \frac{\rho}{\epsilon_0}
                    \label{eq:maxwell1}
                \end{equation}
            \end{block}
        \end{column}
    \end{columns}
\end{frame}
\begin{frame}[fragile]
    \frametitle{Die \texttt{align}-Umgebung \hfill\doc{http://ftp.fau.de/ctan/macros/latex/required/amslatex/math/amsldoc.pdf}{amsmath}}
    \begin{columns}[t]
        \begin{column}{0.35\textwidth}
            \begin{itemize}
                \item für mehrere Gleichungen
                \item \texttt{\&} steuert Ausrichtung
            \end{itemize}
        \end{column}
        \begin{column}{0.6\textwidth}
            \begin{itemize}
                \item \texttt{\textbackslash\textbackslash} erzeugt neue Zeile
                \item jede Zeile bekommt eine Gleichungsnummer
            \end{itemize}
        \end{column}
    \end{columns}
    \vfill
    \begin{columns}[T]
        \begin{column}{0.65\textwidth}
            \begin{block}{\LaTeX-Code}
                \begin{lstverbatim}
                \begin{align}
                    a &= 1       & b&=2 \\
                    a\cdot b &=5 & \frac{a}{b} &= \num{0,5}
                \end{align}
                \end{lstverbatim}
            \end{block}
        \end{column}
        \begin{column}{0.3\textwidth}
            \begin{block}{Ergebnis}
                \begin{align}
                    a &= 2       & b&=2 \\
                    a\cdot b &=2 & \frac{a}{b} &= \num{0,5}
                \end{align}
            \end{block}
        \end{column}
    \end{columns}
\end{frame}
\begin{frame}[fragile]
    \frametitle{Symbol-Sammlung}
    Praktischer Link: \\
    \href{http://detexify.kirelabs.org/classify.html}{http://detexify.kirelabs.org/classify.html} \\
    (Symbol malen und LaTeX-Code angezeigt bekommen)
    \begin{columns}[T]
        \begin{column}{0.7\textwidth}
            \begin{block}{\LaTeX-Code}
                \begin{lstverbatim}
                \begin{align}
                    \leq \geq \gg \ll \approx \propto \\
                    \cdot \times \partial \bar{x} \vec{x} \\
                    \pm \mp \infty \partial \nabla \\
                    \int\limits_0^1 \sum\limits_{i=1}^N \\
                    \iint \iiint \oint
                \end{align}
                \end{lstverbatim}
            \end{block}
        \end{column}
        \begin{column}{0.25\textwidth}
            \begin{block}{Ergebnis}
                \begin{align*}
                    \leq \geq \gg \ll \approx \propto \\
                    \cdot \times \bar{x} \vec{x} \\
                    \pm \mp \infty \partial \nabla \\
                    \int\limits_0^1 \sum\limits_{i=1}^N \\
                    \iint \iiint \oint
                \end{align*}
            \end{block}
        \end{column}
    \end{columns}
\end{frame}
\begin{frame}[fragile]
    \frametitle{Konventionen: Variablen, Zahlen, Einheiten, Indizes}
    \begin{itemize}
        \item Variablen werden kursiv gesetzt. Dies geschieht im Mathematikmodus automatisch.
        \item Einheiten werden aufrecht gesetzt und haben ein kleines Leerzeichen (\verb+\,+) Abstand zu ihrer Zahl. Am besten benutzt man hierfür immer \texttt{siunitx}.
        \item Die eulersche Zahl e, das imaginäre i und das infinitesimale d werden ebenfalls aufrecht gesetzt. Im Mathematikmodus erreicht man dies mit 
            \begin{lstverbatim}
            \mathrm{e}, \mathrm{d}, \mathrm{i}.
            \end{lstverbatim}
        \item Bestehen Indizes aus Text, wie min oder max, so wird dies ebenfalls aufrecht gesetzt.
            \begin{lstverbatim}
            x_\text{min}
            \end{lstverbatim}
        \item ein $\mathrm{d}x$ sollte durch ein kleines Leerzeichen (\verb+\,+) vom Integranden abgetrennt werden.
    \end{itemize}
\end{frame}
\begin{frame}
    \frametitle{Übung: Maxwell-Gleichungen}
    Erstellt mit Hilfe der \texttt{align}-Umgebung die Maxwellgleichungen:
    \begin{block}{Ergebnis}
        \begin{align}
            \nabla \cdot \vec{E} &= \frac{\rho} {\epsilon_0} &
            \nabla \cdot \vec{B} &= 0 \\
            \nabla \times \vec{E} &= - \partial_t \vec{B} &
            \nabla \times \vec{B} &= \mu_0 \vec{j} + \mu_0 \epsilon_0 \partial_t \vec{E}
            \label{align}
        \end{align}
    \end{block}
\end{frame}

\section{Floating-Umgebungen}
\begin{frame}
    \frametitle{Gleitumgebungen}
    Zum setzen Elementen, die nicht zum Fließtext gehören, werden Gleitumgebungen genutzt. Diese werden automatisch an eine passende Stelle gesetzt.
    \begin{itemize}
        \item für Abbildungen und Tabellen
        \item Die Freiheit die \LaTeX \ beim setzen hat, kann mit optionalen Argumenten gestuert werden. 
        \item h  (here), !h (noch strenger), t (top), b (bottom)
    \end{itemize}
\end{frame}
\begin{frame}[fragile]
    \frametitle{Bilder einbinden}
    \begin{block}{benötigte Pakete}
        \begin{lstverbatim}
        \usepackage{graphicx}
        \usepackage[labelfont=bf]{caption}
        \end{lstverbatim}
    \end{block}
    \begin{columns}[T]
        \begin{column}{0.6\textwidth}
            \begin{block}{\LaTeX-Code}
                \begin{lstverbatim}
                \begin{figure}
                    \centering
                    \includegraphics[width=\textwidth]{./peplogo.pdf}
                    \caption{Das Pep-Logo}
                    \label{fig:peplogo}
                \end{figure}
                \end{lstverbatim}
            \end{block}
        \end{column}
        \begin{column}{0.35\textwidth}
            \begin{block}{Ergebnis}
                \begin{figure}
                    \centering
                    \includegraphics[width=\textwidth]{./peplogo.pdf}
                    \caption{Das PeP-Logo}
                    \label{fig:peplogo}
                \end{figure}
            \end{block}
        \end{column}
    \end{columns}
\end{frame}



\end{document}
